% !TEX program = xelatex
% v30: generated from actual saved PNG evaluation reports; self-contained source.
\documentclass[10pt,aspectratio=169]{beamer}
\usetheme{Madrid}
\usecolortheme{dolphin}
\usepackage[fontset=none]{ctex}
% Use two TeX-distributed fonts explicitly; no machine-specific Noto dependency.
\setCJKmainfont{FandolSong-Regular.otf}[BoldFont=FandolSong-Bold.otf,AutoFakeSlant=.2]
\setCJKsansfont{FandolSong-Regular.otf}[BoldFont=FandolSong-Bold.otf,AutoFakeSlant=.2]
\setCJKmonofont{FandolSong-Regular.otf}
\usepackage{amsmath,amssymb,bm,mathtools}
\usepackage{booktabs,tabularx,array,multirow}
\usepackage{graphicx,hyperref}
\usefonttheme{professionalfonts}
\setbeamertemplate{navigation symbols}{}
\setbeamertemplate{footline}[frame number]
\setbeamertemplate{equation number}[default]
\setbeamerfont{frametitle}{size=\large}
\setbeamersize{text margin left=7mm,text margin right=7mm}
\setlength{\parskip}{2pt}
\definecolor{deepblue}{RGB}{27,67,121}
\definecolor{softblue}{RGB}{236,243,252}
\definecolor{softgreen}{RGB}{235,247,241}
\definecolor{softorange}{RGB}{252,243,233}
\definecolor{deepgreen}{RGB}{32,105,80}
\definecolor{deeporange}{RGB}{174,87,20}
\setbeamercolor{block title}{fg=white,bg=deepblue}
\setbeamercolor{block body}{bg=softblue,fg=black}
\setbeamercolor{block title example}{fg=white,bg=deepgreen}
\setbeamercolor{block body example}{bg=softgreen,fg=black}
\setbeamercolor{block title alerted}{fg=white,bg=deeporange}
\setbeamercolor{block body alerted}{bg=softorange,fg=black}
\setbeamercolor{alerted text}{fg=deeporange}
\newcommand{\E}{\mathbb E}
\newcommand{\R}{\mathbb R}
\newcommand{\N}{\mathcal N}
\newcommand{\sg}{\operatorname{sg}}
\newcommand{\Cov}{\operatorname{Cov}}
\newcommand{\KL}{\operatorname{KL}}
\newcommand{\ESS}{\operatorname{ESS}}
\newcommand{\dd}{\,\mathrm d}
\newcommand{\takeaway}[1]{\begin{exampleblock}{本页结论}\small #1\end{exampleblock}}
\title[方案 A 修订：PC-MH]{方案 A 修订：教师引导后验 MH\\与后验去噪一致性}
\subtitle{RTDMD 条件均值估计、可验证的局部校正与单卡实验\\v30：基于 v29 的理论修订与实验更新}
\author{组会讲解版}
\date{2026 年 10 月 21 日}
\hypersetup{pdftitle={RTDMD v30 posterior MH and posterior consistency},pdfauthor={YYPF}}
\begin{document}
\maketitle

\section{目标、保留内容与需要修订的假设}
\begin{frame}{v30 修改了什么，保留了什么}
\small
\begin{columns}[T,onlytextwidth]
\column{.49\textwidth}
\begin{block}{保留 v29 的研究目标}
\begin{itemize}
\item 当前 Student 定义 target；教师参与 proposal 引导。
\item 在线估计条件均值，通过 Tweedie 关系获得 score。
\item 不额外训练 Fake；使用 score gap 的 surrogate 更新。
\item 推理仍为原 CPS 四步，不添加终端噪声。
\end{itemize}
\end{block}
\column{.49\textwidth}
\begin{alertblock}{明确改写的部分}
\begin{itemize}
\item 在真实生成随机变量空间定义 Bayes 后验。
\item 全路径 IS 改为分块 MH；不使用奇异末步密度比。
\item 明确同一查询点的 Student Jacobian。
\item 新增去噪一致性正则，并报告主/敏感性两组实验。
\end{itemize}
\end{alertblock}
\end{columns}
\takeaway{v30 是方法修订版；不是对原 Pref 反向 CPS 全路径 IS 的严格复现。}
\end{frame}

\begin{frame}{RTDMD 的查询与目标 score}
\small
当前 Student 产生源 latent $X_{\rm src}$，独立加噪：
\begin{equation}\label{eq:query}
y=\alpha X_{\rm src}+\sigma\varepsilon,\qquad
\alpha=1-\sigma,\quad \varepsilon\sim\N(0,I),\quad 0<\sigma<1.
\end{equation}
在固定条件 $c$ 下，所需条件均值与 score 为
\begin{equation}\label{eq:tweedie}
m_\theta(y)=\E_\theta[X_{\rm src}\mid y,c],\qquad
s_{\theta,\sigma}(y)=\frac{\alpha m_\theta(y)-y}{\sigma^2}.
\end{equation}
\begin{block}{原 Fake 的作用}
Fake 通过当前 Student 样本的去噪回归近似 $m_\theta(y)$ 或相应 score。
省去 Fake 后，仍然必须估计\alert{这个前向加噪模型的条件均值}。
\end{block}
\takeaway{Tweedie 恒等式保留；本轮改变的是条件均值的构造与训练 surrogate。}
\end{frame}

\begin{frame}{源 latent 必须与实际训练采样器一致}
\small
\begin{equation}\label{eq:source-map}
U=(Z,\xi_1,\ldots,\xi_{K-1},J)\sim P,\qquad
X_{\rm src}=F_\theta(U,c).
\end{equation}
\begin{itemize}\itemsep4pt
\item $Z,\xi_j$ 独立标准 Gaussian；$J$ 按原 DMD 的随机停止分布采样。
\item $F_\theta$ 执行相同 SD3、sigma schedule、CPS 与单条件 Student 预测。
\item 返回第 $J$ 个执行步的 \textbf{clean prediction}，不是噪声边界 $x_s$。
\item 包括未使用的 $\xi_j$，使原始随机变量空间维数固定；它们被积分掉。
\end{itemize}
\begin{alertblock}{有限步训练分布与完整推理分布}
训练源是随机停止的 clean prediction 混合分布；最终评估仍执行完整四步。
两者采用原 DMD 的约定，但不能在理论中混用。
\end{alertblock}
\end{frame}

\begin{frame}{v29 的两个数学缺口}
\small
\begin{block}{一：CPS 的零方差末步}
CPS 创新标准差 $\nu_k=s_{k-1}\sin(\eta\pi/2)$。当 $s_{k-1}=0$，
$M^a_k(\dd x'\mid y)=\delta_{D_a(y)}(\dd x')$。
若两模型输出不同，给定 $y$ 的终端条件核可互相奇异，不能继续代入 Gaussian 密度比。
\end{block}
\begin{block}{二：反向路径终点不自动等于前向后验}
路径比可恢复\emph{所定义的反向 Student 路径律}；仍需另证它的终点条件律等于
式\eqref{eq:query} 的 Bayes 后验。有限步网络本身不提供该保证。
\end{block}
\takeaway{本版将 $F_\theta$ 看作生成映射，直接定义前向后验；不补终端噪声，也不省略末步后伪称原目标。}
\end{frame}

\begin{frame}{原路径比恒等式仍成立，但结论应附带条件}
\small
对正方差的可计算反向核，设
\begin{equation}
R_a(\dd y_{0:K-1}\mid y_K)=\prod_{k=1}^{K}M_k^a(\dd y_{k-1}\mid y_k),\qquad
\Lambda=\frac{\dd R_\theta}{\dd R_{\rm ref}}.
\end{equation}
若绝对连续、可积，则
\begin{equation}
\E_{R_{\rm ref}}[\Lambda Y_0\mid y_K]
=\E_{R_\theta}[Y_0\mid y_K].
\end{equation}
\begin{alertblock}{需要补充的关键条件}
右侧等于 $m_\theta(y_K)$，还要求 $R_\theta(Y_0\mid y_K)$ 正好是
\emph{实际源分布与指定 Gaussian 前向观测}的反向条件律。
\end{alertblock}
\takeaway{原桥分布消元和路径比推导可作为条件性背景；不是本版后验正确性的证明。}
\end{frame}

\section{真实源后验与教师引导局部校正}
\begin{frame}{在原始随机轨迹空间定义真正的 Bayes 后验}
\small
\begin{equation}\label{eq:posterior}
\Pi_\theta(\dd U\mid y,c)
=\frac{P(\dd U)L_\sigma(y\mid F_\theta(U,c))}
{\int L_\sigma(y\mid F_\theta(V,c))P(\dd V)},
\end{equation}
其中对固定 $y,\sigma$，与 $U$ 无关的 Gaussian 常数可消去：
\begin{equation}\label{eq:likelihood}
\ell(U)\equiv\log L_\sigma(y\mid F_\theta(U,c))
=-\frac{\|y-\alpha F_\theta(U,c)\|^2}{2\sigma^2}.
\end{equation}
\begin{equation}\label{eq:true-mean}
m_\theta(y)=\int F_\theta(U,c)\Pi_\theta(\dd U\mid y,c).
\end{equation}
\takeaway{$F_\theta$ 无需可逆，也无需有图像空间密度；确定性末步没有额外核比。}
\end{frame}

\begin{frame}{为什么无条件 Source-IS 在高维退化}
\small
上一版从 $Q_y$ 采样 $U_i$，使用
\begin{equation}
\ell_i=-\frac{\|y-\alpha F_\theta(U_i)\|^2}{2\sigma^2}
+\log\frac{\dd P}{\dd Q_y}(U_i),\qquad
\widetilde w_i=\operatorname{softmax}_i(\ell_i).
\end{equation}
\begin{itemize}\itemsep3pt
\item 源 latent 为 $16\times64\times64=65536$ 维；误差必须逐坐标求和。
\item 8 个近先验候选往往不能覆盖后验；小 $\sigma$ 还会放大失配。
\item 教师估计的腐化噪声不保证等于 Student 的生成初始噪声。
\end{itemize}
\begin{alertblock}{已完成旧运行的定量证据}
ESS 均值 $1.095/8$；96\% 更新 ESS$<2$；log-weight 跨度中位数 $1706.605$。
减去最大 log-weight 只能稳定计算，不能恢复缺失的后验覆盖。
\end{alertblock}
\end{frame}

\begin{frame}{成对源轨迹提供后验平稳初始化}
\small
每轮固定当前 $\theta$，重新生成
\begin{equation}
U^\star\sim P,\qquad X^\star=F_\theta(U^\star,c),\qquad
Y=\alpha X^\star+\sigma\varepsilon.
\end{equation}
联合律为 $P(\dd U)L_\sigma(\dd y\mid F_\theta(U,c))$。由 Bayes 公式，
\begin{equation}\label{eq:paired-stationary}
\mathcal L(U^\star\mid Y=y,c,\theta)=\Pi_\theta(\cdot\mid y,c).
\end{equation}
\begin{block}{成立条件}
必须用\alert{同一次生成的真实轨迹}和其 Gaussian 观测；参数更新后重新生成。
教师图像、旧参数轨迹或独立的另一条源轨迹不能直接替代。
\end{block}
\takeaway{一个条件后验样本不等于精确条件均值；平稳初始化不等于链样本相互独立。}
\end{frame}

\begin{frame}{教师引导的 Gaussian-reference 分块提议}
\small
教师仍在相同 $y,\sigma,c$ 上以原 CFG 预测 $v_{\rm ref}$，固定
\begin{equation}
h_y=\kappa\{y+(1-\sigma)v_{\rm ref}(y,\sigma,c)\},\qquad\kappa=.01.
\end{equation}
固定 $J$，选择一个实际使用的 Gaussian 变量块，再选坐标集合 $C$：
\begin{equation}\label{eq:pcn}
u'_C=\rho u_C+(1-\rho)h_C+\beta\zeta,\qquad
\rho=\sqrt{1-\beta^2},\quad\zeta\sim\N(0,I).
\end{equation}
初始噪声 $Z$ 使用 $h_C=(h_y)_C$；CPS 创新使用 $h_C=0$。
其他坐标保持不变。
\begin{alertblock}{教师中心的含义}
这是可校正的启发式提议中心，不是已证明的生成噪声反演，也不保证优于零中心提议。
\end{alertblock}
\end{frame}

\begin{frame}{接受比必须包含教师中心的校正项}
\small
提议对 $\N(h_C,I)$ 可逆，而真实先验对 $\N(0,I)$ 定义。
故正反向提议比与先验比合并为
\begin{equation}\label{eq:mh-accept}
A(U,U')=\ell(U')-\ell(U)-h_y^\top(Z'-Z),\qquad
a(U,U')=\min\{1,e^{A(U,U')}\}.
\end{equation}
\begin{itemize}\itemsep3pt
\item 未变化的坐标对最后一项贡献为零；零中心创新块没有额外中心项。
\item 活跃块的选择概率只依赖\alert{未改变的} $J$，正反向相同。
\item $J$ 单独用均匀先验提议，保持 Gaussian 变量不变，接受比仅为似然差。
\end{itemize}
\takeaway{每个转移保持式\eqref{eq:posterior} 的后验不变；不能只按观测似然接受一个非零中心提议。}
\end{frame}

\begin{frame}{当前局部核：保持目标，也留下混合局限}
\small
\begin{equation}
\beta=\operatorname{clip}\left(
\frac{.25\sigma}{(1-\sigma)\sqrt{16}}\,b,\;10^{-4},\;.5\right),
\qquad b\text{ 循环取 }(.5,1,2).
\end{equation}
\begin{itemize}\itemsep3pt
\item 每次更新 8 个提议：7 次 Gaussian 分块更新，最后 1 次停止步提议。
\item 块大小16；噪声相关尺度只依赖固定 $\sigma$，不按本链接受历史随意自适应。
\item 最多提议 $7\times16=112$ 个连续坐标（可重复），远小于完整随机轨迹维数。
\item BF16 网络、小噪声尖锐似然与局部更新可能共同限制移动；未单独做精度消融。
\end{itemize}
\takeaway{后验不变性与混合效率是两个问题；本轮正则化也不被宣称为解决链混合。}
\end{frame}

\begin{frame}{Rao--Blackwell 化接受／拒绝，得到在线均值}
\small
当前状态 $U_j$、提议 $U'_j$ 和接受概率 $a_j$ 给出
\begin{equation}
H_j=(1-a_j)F_\theta(U_j,c)+a_jF_\theta(U'_j,c),\qquad
\widehat m=\frac1M\sum_{j=0}^{M-1}H_j.
\end{equation}
$H_j$ 是对接受随机数积分后的下一状态条件期望。
由平稳初始化及不变核，理想连续算术模型下，
\begin{equation}\label{eq:conditional-mean}
\E[\widehat m\mid y,c,\sigma,\theta]=m_\theta(y).
\end{equation}
\begin{alertblock}{这条结论不包含什么}
不表示一次估计精确，不表示8步已充分混合；浮点计算仍有数值近似。
对任意非平稳初始化，也不能直接沿用式\eqref{eq:conditional-mean}。
\end{alertblock}
\end{frame}

\section{同一观测点的更新与新增一致性正则}
\begin{frame}{为什么要明确 Jacobian 的位置}
\small
上一版沿用生成源样本 Jacobian。成对轨迹均值估计与该源样本条件相关，
一般不能将它们的期望乘开。例如 $X\mid y$ 的均值为 $m$、方差为 $v$，则
\begin{equation}
\E[X\cdot X\mid y]=m^2+v\ne m^2.
\end{equation}
本版按 v29 式(2)的一种明确解释，在同一查询点使用
\begin{equation}\label{eq:denoiser}
D_\theta(y)=y-\sigma v_\theta(y,\sigma,c),\qquad J_\theta D_\theta(y).
\end{equation}
给定 $y,c,\sigma,\theta$，这个 Jacobian 是确定的；查询、轨迹和目标均停止梯度。
实现要求 Student 无正概率 Dropout。
\takeaway{这不是原 DMD 源 Jacobian 的等价改写；实验收益尚需单独消融两种改动。}
\end{frame}

\begin{frame}{Posterior-MH 的基础 surrogate 与时间预条件}
\small
令教师 clean prediction $r=y-\sigma v_{\rm ref}$，定义
\begin{equation}
n_\theta(y)=\operatorname{mean}|D_\theta(y)-r|>0,\qquad
G_\theta(y)=\frac{1}{d\,n_\theta(y)}\{J_\theta D_\theta(y)\}^\top.
\end{equation}
保留仓库 DMD 的 clean 差归一化形式，基础更新为
\begin{equation}\label{eq:base-surrogate}
\widehat g_0=G_\theta(y)(\widehat m-r).
\end{equation}
$r$ 是冻结教师的 CFG 网络预测，不是已经证明的精确后验均值；
以下代数命题只要求给定观测后 $r$ 确定，不声称 CFG 目标的 score 精确。
\begin{block}{与 Tweedie score 差的关系}
真实 score 差还含 $\alpha/\sigma^2$。上述更新将其吸收于时间权重，
另对状态相关的 $n_\theta$ 停止梯度；不是未指定时间权重下的等价 KL 梯度。
\end{block}
\takeaway{以下结论针对固定 $\theta,y$ 的裁剪前 surrogate 梯度，不是完整 KL 梯度或收敛定理。}
\end{frame}

\begin{frame}{v30 新增：后验去噪一致性}
\small
网络的 clean prediction 与其\emph{实际源分布}的条件均值不一定一致。
在固定目标和归一化下，加入半梯度正则
\begin{equation}\label{eq:consistency}
\mathcal L_{\rm pc}=\frac{\|D_\theta(y)-\sg(\widehat m)\|^2}
{2d\,\sg(n_\theta(y))}.
\end{equation}
其条件期望梯度为 $G_\theta(y)\{D_\theta(y)-m_\theta(y)\}$。
和式\eqref{eq:base-surrogate} 合并得到
\begin{equation}\label{eq:regularized}
\widehat g_\lambda
=G_\theta(y)\left\{(1-\lambda)\widehat m+\lambda D_\theta(y)-r\right\}.
\end{equation}
\takeaway{这是新增的 Student 后验一致性项，不是 v29 的 Fake 一致性正则 $\mathcal L_{\rm cons}$。
本节 $\lambda$ 是正则强度，不是原稿的时间权重 $\lambda_t$。}
\end{frame}

\begin{frame}{新增命题：条件后验采样方差的精确缩放}
\small
\begin{block}{条件与结论}
固定 $y,c,\sigma,\theta$；$G_\theta,D_\theta,r$ 确定，$\lambda$ 预先固定，
后验均值估计二阶矩有限，归一化为正。则裁剪前
\begin{equation}\label{eq:variance}
\Cov(\widehat g_\lambda\mid y)
=(1-\lambda)^2\Cov(\widehat g_0\mid y).
\end{equation}
\end{block}
\begin{proof}[一行证明]
将式\eqref{eq:regularized} 写成
$(1-\lambda)G_\theta\widehat m+G_\theta(\lambda D_\theta-r)$；
第二项给定观测后确定，取协方差即得。
\end{proof}
$\lambda=.25$ 时系数为 $.5625$；$\lambda=.50$ 时为 $.25$。
\begin{alertblock}{严格限定}
只缩放\alert{给定观测的后验估计噪声}；不等于总 SGD 方差、Adam 更新方差或图像指标的同幅改善。
\end{alertblock}
\end{frame}

\begin{frame}{期望更新也改变了：不是无代价的方差缩减}
\small
由式\eqref{eq:conditional-mean}，
\begin{equation}\label{eq:expected-regularized}
\E[\widehat g_\lambda\mid y]
=G_\theta(y)\{m_\theta(y)-r\}
+\lambda G_\theta(y)\{D_\theta(y)-m_\theta(y)\}.
\end{equation}
\begin{itemize}\itemsep4pt
\item 新增项正是去噪一致性修正；因此 $\lambda>0$ 是正则化扩展。
\item 若同时 $m_\theta=r$、$D_\theta=m_\theta$，该目标点的两个项均为零。
\item 不由此推出解唯一、参数稳定、KL 单调下降或全局收敛。
\item 源分布随 $\theta$ 变化；本实现对采样数据和目标使用 stop-gradient。
\end{itemize}
\takeaway{方差命题与期望更新偏移必须同时报告，不能把它表述为保持原方法不变的无偏改进。}
\end{frame}

\begin{frame}{边界情形与本轮预先固定的实验}
\small
\begin{center}
\begin{tabular}{cll}\toprule
$\lambda$ & 更新中的向量 & 含义\\\midrule
$0$ & $\widehat m-r$ & 上一版 Posterior-MH\\
$.25$ & $.75\widehat m+.25D_\theta-r$ & 本轮主实验\\
$.50$ & $.50\widehat m+.50D_\theta-r$ & 本轮敏感性实验\\
$1$ & $D_\theta-r$ & 理论边界：后验项完全抵消\\\bottomrule
\end{tabular}
\end{center}
\begin{alertblock}{实验选择规则}
主实验 $.25$ 和敏感性 $.50$ 在最终测试评分前固定；两组均报告。
不按最终2048条测试分数挑选一个强度，隐去另一组结果。
\end{alertblock}
\takeaway{$\lambda=1$ 不再保留后验均值的作用；不能把纯教师拟合的效果当成 MH 的贡献。}
\end{frame}

\begin{frame}{训练算法与实际实现的对应}
\small
\begin{enumerate}\itemsep4pt
\item 用当前 $\theta$ 采样真实轨迹 $U^\star$，重放得到 $X^\star$；Gaussian 加噪得到 $y$。
\item 教师在 $y$ 预测 $r,h_y$；固定 $\theta$，执行8次带校正的局部 MH，得到 $\widehat m$。
\item Student 在\alert{同一固定 $y$}计算 $D_\theta(y)$；构造式\eqref{eq:regularized} 的停止梯度目标。
\item 仅通过 $D_\theta(y)$ 反传；使用原学习率、优化器、梯度裁剪和 EMA。
\item 原 DMD 每10外层轮一次生成器更新的日程保留；不执行 Fake 更新。
\end{enumerate}
\begin{block}{额外计算}
一致性项通过向量合并实现，不增加网络前向次数。MH 候选仍然需要额外回放；
训练计算量与 DMD 不相等，推理 NFE 保持4。
\end{block}
\end{frame}

\section{实验协议、结果与诊断}
\begin{frame}{公平比较：保持的参数与必要差异}
\small
\begin{columns}[T,onlytextwidth]
\column{.49\textwidth}
\begin{block}{共同设置}
原 SD3-Medium；512；BF16 网络；LoRA rank32/alpha64；batch1；seed42；
学习率 $3\times10^{-5}$；相同优化器、EMA、裁剪、CFG4.5；CPS $\eta=.9$；
四步 $[1000,750,500,250]$；随机停止；相同提示词与评估种子。
\end{block}
\column{.49\textwidth}
\begin{alertblock}{必须披露的差异}
无 Fake 训练；同点 Jacobian；8次 MH、16坐标块及一致性系数。
源/观测 Gaussian 先以FP32采样，旧DMD直接BF16采样。
只保存末轮 EMA LoRA，未保存完整可恢复训练状态。
\end{alertblock}
\end{columns}
\takeaway{相同外部配置与生成器更新数，不代表训练计算量相等；本轮没有启动16卡任务。}
\end{frame}

\begin{frame}{数据、评分与本轮资源使用}
\small
\begin{itemize}\itemsep5pt
\item PickScore 训练提示词 25432 条；测试提示词 2048 条，原顺序、每条一张图。
\item 1000外层轮，100次生成器更新；这不是默认长训练或收敛复现。
\item 四项指标评分\alert{相同保存PNG}：CLIP、Aesthetic、PickScore、HPSv2.1。
\item CLIP 与 HP 展示值乘100；Pick还原原始logit，Aesthetic不缩放。
\item 两个新增实验共享同一A100并行运行；因此不比较它们与旧运行的墙钟速度。
\item 完整图片暂存内存盘；权重、逐图分数、配置、报告保存在项目目录。
\end{itemize}
\takeaway{测试集用于最终报告；本轮正则强度预先固定。跨运行多次查看测试分数仍限制外部泛化结论。}
\end{frame}

\begin{frame}{完整2048条测试提示词：全部方法结果}
\small
\begin{center}
\setlength{\tabcolsep}{4pt}
\begin{tabular}{lrrrrr}\toprule
Post-Train（均无奖励） & NFE & CLIP$\uparrow$ & Aesthetic$\uparrow$ & Pick$\uparrow$ & HP$\uparrow$\\\midrule
DMD & 4 & 25.5374 & 5.9152 & 20.3478 & 27.2742 \\
Source-IS & 4 & 19.5789 & 4.5526 & 17.9098 & 16.0197 \\
Posterior-MH & 4 & 25.4264 & 5.9660 & 20.5580 & 28.3938 \\
PC-MH, $\lambda=.25$ & 4 & 26.1713 & 6.0386 & 20.8326 & 28.0147 \\
PC-MH, $\lambda=.50$ & 4 & 26.2382 & 5.9494 & 20.7121 & 26.5751 \\
\bottomrule
\end{tabular}
\end{center}
\begin{block}{预先指定的主实验比较}
预先指定的主实验相对 DMD：CLIP +0.6339；Aesthetic +0.1234；Pick +0.4848；HP +0.7405。\par 相对上一版 Posterior-MH：CLIP +0.7448；Aesthetic +0.0726；Pick +0.2746；HP -0.3790。
\end{block}
\begin{alertblock}{结论边界}
一个训练种子、100次生成器更新。指标更高也不构成稳定优越性、理论正确性或充分收敛的证明。
\end{alertblock}
\end{frame}

\begin{frame}{本轮结果怎样支持下一步，而不是过度归因}
\small
\begin{itemize}\itemsep6pt
\item 主实验 $\lambda=.25$ 相对 DMD 四项提高；相对上一版 MH，CLIP、Aesthetic、Pick 提高，但 HP 下降。
\item 敏感性 $\lambda=.50$ 进一步提高 CLIP，但 HP 低于 DMD；说明正则强度存在指标之间的权衡。
\item 选择规则保持不变：$.25$ 是预先指定的主实验；$.50$ 是敏感性实验，不改成事后最优方法。
\item 这些观测支持继续检验一致性正则，但不能证明收益来自条件方差缩减，更不能证明 MH 的独立贡献。
\end{itemize}
\takeaway{先补独立验证集、多训练种子及0步MH消融；当前单次短跑不支持直接推断16卡长训练的收益。}
\end{frame}

\begin{frame}{链诊断：接受不等于实际移动，更不等于ESS}
\small
\begin{center}
\begin{tabular}{lrrr}\toprule
方法 & 平均接受率 & 实际图像移动率 & 无实际移动的更新比例\\\midrule
Posterior-MH & 28.75\% & 25.75\% & 50.0\% \\
PC-MH, $\lambda=.25$ & 33.50\% & 30.50\% & 44.0\% \\
PC-MH, $\lambda=.50$ & 36.75\% & 33.75\% & 40.0\% \\
\bottomrule
\end{tabular}
\end{center}
\begin{itemize}\itemsep4pt
\item 上一版低于 $.5$ 的噪声段没有实际图像移动；全体50\%更新无图像移动。
\item BF16舍入可能使原始噪声改变而输出图像不变；接受率会高于实际移动率。
\item 一致性正则不改变 MH 的目标或核，不应被宣称为修复混合。
\item 8次转移不足以可靠估计MCMC ESS；不把等权平均报告为 ESS$=8$。
\end{itemize}
\end{frame}

\begin{frame}{数学验证与代码证据}
\small
\begin{block}{20项测试通过}
原6项：Gaussian密度比、解析后验/Tweedie、非可逆末步、维度求和等；
MH8项：Hastings比、平稳性、随机停止、轨迹重放、相关Gaussian块、冻结诊断等；
新增6项：$\lambda=0/1$边界、方差系数、期望方向改变、目标停止梯度、参数范围。
\end{block}
\begin{itemize}\itemsep4pt
\item 实际代码、数据SHA256与配置均保存；旧DMD和上一版结果未覆盖。
\item 完整逐图评分与提示词索引可核验，不只保存均值。
\item 测试验证实现中的指定关系，不替代神经网络全局训练收敛分析。
\end{itemize}
\takeaway{论文数值表从实际评分JSON生成；缺失结果不会以估计数字填充。}
\end{frame}

\begin{frame}{当前仍缺少的消融与统计证据}
\small
\begin{enumerate}\itemsep5pt
\item 0步MH：直接使用成对源样本，判断MH在同点Jacobian下的增量作用。
\item 教师偏移设为0：判断当前 $\kappa=.01$ 是否改善采样而非仅由教师目标起作用。
\item 分别改变后验估计、Jacobian及一致性项，避免把所有收益归给一个因素。
\item 多个训练种子；独立验证提示词；等更新数和GPU小时两种成本比较。
\item 对低噪声的BF16回放、块大小/步长、长链相关性进行独立诊断。
\end{enumerate}
\takeaway{本轮为有理论定义的实验修订；没有完成这些消融，因此不作完整因果归因。}
\end{frame}

\begin{frame}{扩大到16卡之前应完成什么}
\small
\begin{itemize}\itemsep5pt
\item 优先将可用GPU用于独立种子与关键消融，避免只放大单次试跑。
\item 先验证2--4卡同步、每卡独立轨迹随机数、梯度平均和EMA一致性。
\item 补齐原始参数、优化器、调度器、EMA、随机数和数据位置的完整断点恢复。
\item 正式16卡对照需匹配全局batch、生成器更新数、数据暴露量，并报告GPU小时。
\end{itemize}
\begin{alertblock}{当前代码能力}
当前回放入口明确限制单卡batch1；导出的EMA LoRA不是完整训练恢复。
16卡数据并行也不会自动改善每个观测的后验链混合。
\end{alertblock}
\end{frame}

\appendix
\section{附录：原CPS与路径比背景的适用边界}
\begin{frame}{附录 A：原CPS生成映射保持不变}
\small
对RF速度 $v_\theta(y,s)$，令
\begin{equation}
D_\theta(y,s)=y-sv_\theta(y,s),\qquad
\widehat\varepsilon=y+(1-s)v_\theta(y,s).
\end{equation}
CPS 从 $s$ 到 $t<s$：
\begin{equation}
x_t=(1-t)D_\theta(y,s)+t\cos(\eta\pi/2)\widehat\varepsilon
+t\sin(\eta\pi/2)\xi.
\end{equation}
$t=0$ 时得到确定性clean输出。该映射可以直接重放；不需要把最后一步当成有密度的Gaussian核。
\takeaway{v30 不修改推理CPS、噪声强度或NFE；修改发生在在线估计和训练更新。}
\end{frame}

\begin{frame}{附录 B：正方差共同协方差核的比值}
\small
若两种辅助反向核确有共同正定协方差 $\Sigma_k$，且均值为 $\mu_a$，则
\begin{equation}
\log\frac{M_k^\theta(y'\mid y)}{M_k^{\rm ref}(y'\mid y)}
=-\frac12\left[\|y'-\mu_\theta(y)\|_{\Sigma_k^{-1}}^2
-\|y'-\mu_{\rm ref}(y)\|_{\Sigma_k^{-1}}^2\right].
\end{equation}
这是 v29 Gaussian核比与离散Girsanov型展开的代数起点。
\begin{alertblock}{边界}
协方差退化时不能直接套用；即使路径校正精确，也只恢复所定义辅助反向路径律。
它是否对应真正前向后验仍需证明。本版MH不以此式作核心实现。
\end{alertblock}
\end{frame}

\begin{frame}{附录 C：一致性正则的实际surrogate实现}
\small
将所需向量合并为
\begin{equation}
b_\lambda=\sg\!\left[
\frac{(1-\lambda)\widehat m+\lambda D_\theta(y)-r}{n_\theta(y)}\right],
\quad
\widetilde L_\lambda=\frac12\operatorname{MSE}\left(D_\theta(y),\sg[D_\theta(y)-b_\lambda]\right).
\end{equation}
\begin{itemize}\itemsep4pt
\item 目标分支整体停止梯度；只更新Student同一点预测。
\item 原代码的有限性处理、梯度裁剪和Adam继续沿用。
\item 方差命题针对裁剪前；非线性裁剪和优化器之后不能直接套同一系数。
\item $n_\theta=0$等退化点需要数值处理，理想命题假定正归一化。
\end{itemize}
\end{frame}

\begin{frame}{参考与可复现材料}
\footnotesize
\begin{itemize}\itemsep5pt
\item 用户提供的 \texttt{RTDMD\_SchemeA\_Reorganized\_Final\_CN\_v29.tex}：v30的背景、符号和原始方案来源。
\item Cotter, Roberts, Stuart, White. \emph{MCMC Methods for Functions}. Statistical Science, 2013.
\url{https://arxiv.org/abs/1202.0709}
\item Snyder et al. \emph{Obstacles to High-Dimensional Particle Filtering}. Monthly Weather Review, 2008.
\url{https://doi.org/10.1175/2008MWR2529.1}
\item 实验入口：\texttt{train\_rtdmd\_pcmh\_sd3.py}；正则模块：\texttt{posterior\_consistency.py}。
\item 报告：\texttt{logs/rtdmd-pcmh/comparison-v30/}；逐图评分、参数差异、SHA256与存储说明随代码包提供。
\end{itemize}
\end{frame}
\end{document}
